%
% Typography and Publishing - 4th Project
% Ac.Year 2021/2022, Summer Semester
%
% Login xkvace00, Kvacek David
% 1st year BITP, full-time, FIT
% xkvace00@stud.fit.vutbr.cz
%
% Date: 2022-04-16
%

\documentclass[a4paper, 11pt]{article}
\usepackage[left=2cm, text={17cm, 25cm}, top=3cm]{geometry}
\usepackage[czech]{babel}
\usepackage[utf8]{inputenc}
\usepackage[IL2]{fontenc}
\usepackage{hyperref}
\usepackage{times}

\DeclareRobustCommand{\BibTeX}{{\normalfont B\kern-.05em{\scshape i\kern-.025em b}\kern-.08em \TeX}}

\begin{document}

\begin{titlepage}
    \begin{center}
        \Huge\textsc{Vysoké učení technické v~Brně} \\
        \huge\textsc{Fakulta informačních technologií} \\
        \vspace{\stretch{0.382}}
        \LARGE{Typografie a~publikování\,--\,4.~projekt} \\
        \Huge{Bibliografické citace} \\
        \vspace{\stretch{0.618}}
    \end{center}
    \Large{\today \hfill David Kvaček}
\end{titlepage}

\section*{Úvod}
Jakýkoliv dokument, technická zpráva či odborný článek by měl mít svou grafickou kvalitu.
A~právě této oblasti zpracování písemného projevu se věnuje vědní obor typografie.
K~dosažení grafické dokonalosti nám mohou pomoci různé programy, ať už se jedná o~Microsoft Office Word, Apache OpenOffice nebo LibreOffice Writer.
Pravdou ale je, že momentálně nejpokročilejším pomocníkem se zdá být \TeX.
Je možné se setkat i~s~různými nadstavbami jako \LaTeX nebo \BibTeX , který se používá pro bibliografické citace, a~to včetně podpory pro normu ČSN ISO 690 \cite{Martinek2008}.

\section{Písmo}
Základním výstupem našich prací je bezpochyby samotné písmo.
V~minulosti, kdy byl předním nástrojem pro psaní textů psací stroj, se používalo tzv. neproporcionální písmo.
To znamenalo, že šířka všech znaků byla stejná.
Dnes se obvykle používá písmo proporcionální, protože je pro náš zrak příjemnější a~čitelnější.
Písmena se stejnou šířkou se však používají např. pro zvýraznění některé části textu, podrobněji zpracováno v~kvalifikační práci \cite{MUNI2011}.
Navíc je již v~dnešní době podporováno kódování znaků Unicode, které umožňuje šifrovat složitější i~speciální znaky, což využijí především asijské země \cite{Asia2008a}.
S~tím však souvisí přepracování nástroje \BibTeX, na které navazuje článek \cite{Asia2008b}.

\section{Sazba}
Výsledný text může být vysázen hladkou, nebo smíšenou sazbou.
Pro dokumenty odpovídající určité kvalitě je typická hladká sazba.
Ta nepoužívá tolik zdobného textu, aby nenarušila plynulost projevu.
Tomu napomáhá i~automatické dělení slov, které je dále posáno zde \cite{Olsak2014}.
Vyskytují se zde interpunkční znaménka, pomlčka a~spojovník, uvozovky a~závorky a~další znaky (např. apostrof, paragraf, ampersand, hvězdička, procento, promile, stupeň aj.).
Některá další pravidla můžete nalézt v~této práci \cite{VSPJ2021}.

\section{Seznam}
Pokud máme v~úmyslu napsat nějakou posloupnost položek, výčet vlastností nebo seznam osob, můžeme použít speciální prostředí, které je určeno přesně k~tomuto účelu.
\TeX\ navíc umožňuje libovolné nastavení zarovnání odrážek a~jejich vzdálenost.
Tu můžeme zadávat nejen v~imperiálních, ale také pomocí metrických jednotek.
Software je přizpůsoben rovněž na oddělování desetinných čísel čárkou i~tečkou, což je detailněji popsáno v~této monografii \cite{Olsak2016}.

\section{Tabulky}
Nástroj \TeX\ není primárně určen pro sázení tabulek.
Pokud však přesto potřebujeme ve svém textu nějakou pasáž zpřehlednit, tabulka je na to ideální.
S~tímto řešením nám může pomoci nějaké specializované prostředí, které nám ulehčí práci.
Jak pracovat s~jednotlivými nástroji se dozvíte v~\cite{Voss2011}.

\section{Matematika}
Kde je ale \TeX\ doma, to je sázení matematických rovnic, vzorců a~algoritmů.
Podle mého názoru nenajdeme lepší software s~tak propracovaným matematickým prostředím.
Při sázení složitých matematicých výrazů se nemusíme starat o~velikost, stačí se přepnout do určitého módu.
Můžeme si vybrat mezi kompaktním $\sum_{i=1}^n i_n$ a~samostatným režimem \cite{Satrapa2011}.
$$\int_{0}^{\pi} \sqrt{\alpha^{2} + x^{2}}\,dx$$
Tento výraz a~mnoho dalšího naleznete v~\cite{Gratzer2004}.

\newpage
\bibliographystyle{czechiso}
\bibliography{proj4}

\end{document}

